% Options for packages loaded elsewhere
\PassOptionsToPackage{unicode}{hyperref}
\PassOptionsToPackage{hyphens}{url}
\documentclass[
]{article}
\usepackage{xcolor}
\usepackage[margin=1in]{geometry}
\usepackage{amsmath,amssymb}
\setcounter{secnumdepth}{-\maxdimen} % remove section numbering
\usepackage{iftex}
\ifPDFTeX
  \usepackage[T1]{fontenc}
  \usepackage[utf8]{inputenc}
  \usepackage{textcomp} % provide euro and other symbols
\else % if luatex or xetex
  \usepackage{unicode-math} % this also loads fontspec
  \defaultfontfeatures{Scale=MatchLowercase}
  \defaultfontfeatures[\rmfamily]{Ligatures=TeX,Scale=1}
\fi
\usepackage{lmodern}
\ifPDFTeX\else
  % xetex/luatex font selection
\fi
% Use upquote if available, for straight quotes in verbatim environments
\IfFileExists{upquote.sty}{\usepackage{upquote}}{}
\IfFileExists{microtype.sty}{% use microtype if available
  \usepackage[]{microtype}
  \UseMicrotypeSet[protrusion]{basicmath} % disable protrusion for tt fonts
}{}
\makeatletter
\@ifundefined{KOMAClassName}{% if non-KOMA class
  \IfFileExists{parskip.sty}{%
    \usepackage{parskip}
  }{% else
    \setlength{\parindent}{0pt}
    \setlength{\parskip}{6pt plus 2pt minus 1pt}}
}{% if KOMA class
  \KOMAoptions{parskip=half}}
\makeatother
\usepackage{color}
\usepackage{fancyvrb}
\newcommand{\VerbBar}{|}
\newcommand{\VERB}{\Verb[commandchars=\\\{\}]}
\DefineVerbatimEnvironment{Highlighting}{Verbatim}{commandchars=\\\{\}}
% Add ',fontsize=\small' for more characters per line
\usepackage{framed}
\definecolor{shadecolor}{RGB}{248,248,248}
\newenvironment{Shaded}{\begin{snugshade}}{\end{snugshade}}
\newcommand{\AlertTok}[1]{\textcolor[rgb]{0.94,0.16,0.16}{#1}}
\newcommand{\AnnotationTok}[1]{\textcolor[rgb]{0.56,0.35,0.01}{\textbf{\textit{#1}}}}
\newcommand{\AttributeTok}[1]{\textcolor[rgb]{0.13,0.29,0.53}{#1}}
\newcommand{\BaseNTok}[1]{\textcolor[rgb]{0.00,0.00,0.81}{#1}}
\newcommand{\BuiltInTok}[1]{#1}
\newcommand{\CharTok}[1]{\textcolor[rgb]{0.31,0.60,0.02}{#1}}
\newcommand{\CommentTok}[1]{\textcolor[rgb]{0.56,0.35,0.01}{\textit{#1}}}
\newcommand{\CommentVarTok}[1]{\textcolor[rgb]{0.56,0.35,0.01}{\textbf{\textit{#1}}}}
\newcommand{\ConstantTok}[1]{\textcolor[rgb]{0.56,0.35,0.01}{#1}}
\newcommand{\ControlFlowTok}[1]{\textcolor[rgb]{0.13,0.29,0.53}{\textbf{#1}}}
\newcommand{\DataTypeTok}[1]{\textcolor[rgb]{0.13,0.29,0.53}{#1}}
\newcommand{\DecValTok}[1]{\textcolor[rgb]{0.00,0.00,0.81}{#1}}
\newcommand{\DocumentationTok}[1]{\textcolor[rgb]{0.56,0.35,0.01}{\textbf{\textit{#1}}}}
\newcommand{\ErrorTok}[1]{\textcolor[rgb]{0.64,0.00,0.00}{\textbf{#1}}}
\newcommand{\ExtensionTok}[1]{#1}
\newcommand{\FloatTok}[1]{\textcolor[rgb]{0.00,0.00,0.81}{#1}}
\newcommand{\FunctionTok}[1]{\textcolor[rgb]{0.13,0.29,0.53}{\textbf{#1}}}
\newcommand{\ImportTok}[1]{#1}
\newcommand{\InformationTok}[1]{\textcolor[rgb]{0.56,0.35,0.01}{\textbf{\textit{#1}}}}
\newcommand{\KeywordTok}[1]{\textcolor[rgb]{0.13,0.29,0.53}{\textbf{#1}}}
\newcommand{\NormalTok}[1]{#1}
\newcommand{\OperatorTok}[1]{\textcolor[rgb]{0.81,0.36,0.00}{\textbf{#1}}}
\newcommand{\OtherTok}[1]{\textcolor[rgb]{0.56,0.35,0.01}{#1}}
\newcommand{\PreprocessorTok}[1]{\textcolor[rgb]{0.56,0.35,0.01}{\textit{#1}}}
\newcommand{\RegionMarkerTok}[1]{#1}
\newcommand{\SpecialCharTok}[1]{\textcolor[rgb]{0.81,0.36,0.00}{\textbf{#1}}}
\newcommand{\SpecialStringTok}[1]{\textcolor[rgb]{0.31,0.60,0.02}{#1}}
\newcommand{\StringTok}[1]{\textcolor[rgb]{0.31,0.60,0.02}{#1}}
\newcommand{\VariableTok}[1]{\textcolor[rgb]{0.00,0.00,0.00}{#1}}
\newcommand{\VerbatimStringTok}[1]{\textcolor[rgb]{0.31,0.60,0.02}{#1}}
\newcommand{\WarningTok}[1]{\textcolor[rgb]{0.56,0.35,0.01}{\textbf{\textit{#1}}}}
\usepackage{graphicx}
\makeatletter
\newsavebox\pandoc@box
\newcommand*\pandocbounded[1]{% scales image to fit in text height/width
  \sbox\pandoc@box{#1}%
  \Gscale@div\@tempa{\textheight}{\dimexpr\ht\pandoc@box+\dp\pandoc@box\relax}%
  \Gscale@div\@tempb{\linewidth}{\wd\pandoc@box}%
  \ifdim\@tempb\p@<\@tempa\p@\let\@tempa\@tempb\fi% select the smaller of both
  \ifdim\@tempa\p@<\p@\scalebox{\@tempa}{\usebox\pandoc@box}%
  \else\usebox{\pandoc@box}%
  \fi%
}
% Set default figure placement to htbp
\def\fps@figure{htbp}
\makeatother
\setlength{\emergencystretch}{3em} % prevent overfull lines
\providecommand{\tightlist}{%
  \setlength{\itemsep}{0pt}\setlength{\parskip}{0pt}}
\usepackage{bookmark}
\IfFileExists{xurl.sty}{\usepackage{xurl}}{} % add URL line breaks if available
\urlstyle{same}
\hypersetup{
  pdftitle={Posterior Mean and Variance},
  hidelinks,
  pdfcreator={LaTeX via pandoc}}

\title{Posterior Mean and Variance}
\author{}
\date{\vspace{-2.5em}2026-04-17}

\begin{document}
\maketitle

\subsection{Introduction}\label{introduction}

A Bayesian analysis returns a full posterior distribution, but applied
work eventually demands a number to put in a sentence: a coefficient, a
probability, an effect size. The posterior mean, median, and mode are
the three most common point summaries, each motivated by a different
decision-theoretic loss. The posterior variance and standard deviation
quantify the uncertainty around the chosen point. Treating the posterior
as the headline and the point estimate as a one-line summary derived
from it is a more honest reporting style than collapsing the posterior
to a single number first and worrying about uncertainty later.

\subsection{Prerequisites}\label{prerequisites}

A working understanding of Bayes' theorem and the difference between a
posterior distribution and a frequentist sampling distribution; basic
familiarity with MCMC draws or a closed-form posterior.

\subsection{Theory}\label{theory}

For a posterior \(p(\theta \mid y)\), the standard summaries are:

\begin{itemize}
\tightlist
\item
  \textbf{Posterior mean}: \(\mathbb E[\theta \mid y]\) minimises
  expected squared-error loss; it is the Bayes estimator under quadratic
  loss and the natural target of MCMC averaging.
\item
  \textbf{Posterior median}: minimises expected absolute-error loss;
  preferred for skewed posteriors because it does not chase a long tail.
\item
  \textbf{Posterior mode (MAP)}: maximises \(p(\theta \mid y)\) and
  often coincides with a penalised maximum-likelihood estimate when the
  prior is interpreted as a penalty.
\item
  \textbf{Posterior variance}: \(\mathrm{Var}(\theta \mid y)\) measures
  spread; the standard deviation \(\sqrt{\mathrm{Var}}\) is on the
  original scale of \(\theta\).
\end{itemize}

The mean and variance always exist for proper integrable posteriors; the
mode may not be unique under multimodal posteriors and the median is
unique only for continuous distributions.

\subsection{Assumptions}\label{assumptions}

A proper posterior with finite mean and variance, and either closed-form
expressions or sufficient MCMC draws. For heavy-tailed posteriors, the
mean and variance can be unstable estimators even when they exist
mathematically, so report effective sample sizes alongside the
summaries.

\subsection{R Implementation}\label{r-implementation}

\begin{Shaded}
\begin{Highlighting}[]
\FunctionTok{set.seed}\NormalTok{(}\DecValTok{2026}\NormalTok{)}

\CommentTok{\# Simulated samples from a posterior (e.g., from MCMC)}
\NormalTok{theta }\OtherTok{\textless{}{-}} \FunctionTok{rbeta}\NormalTok{(}\DecValTok{5000}\NormalTok{, }\AttributeTok{shape1 =} \DecValTok{10}\NormalTok{, }\AttributeTok{shape2 =} \DecValTok{4}\NormalTok{)}

\FunctionTok{c}\NormalTok{(}\AttributeTok{mean =} \FunctionTok{mean}\NormalTok{(theta),}
  \AttributeTok{median =} \FunctionTok{median}\NormalTok{(theta),}
  \AttributeTok{mode =} \FunctionTok{density}\NormalTok{(theta)}\SpecialCharTok{$}\NormalTok{x[}\FunctionTok{which.max}\NormalTok{(}\FunctionTok{density}\NormalTok{(theta)}\SpecialCharTok{$}\NormalTok{y)],}
  \AttributeTok{sd =} \FunctionTok{sd}\NormalTok{(theta),}
  \AttributeTok{var =} \FunctionTok{var}\NormalTok{(theta))}

\CommentTok{\# Credible interval}
\FunctionTok{quantile}\NormalTok{(theta, }\FunctionTok{c}\NormalTok{(}\FloatTok{0.025}\NormalTok{, }\FloatTok{0.975}\NormalTok{))}
\end{Highlighting}
\end{Shaded}

\subsection{Output \& Results}\label{output-results}

Five scalar summaries and a 95 \% equal-tailed credible interval. For a
Beta(10, 4) posterior the mean is approximately \(10/14 \approx 0.71\),
the standard deviation is around \(0.11\), and the 95 \% interval
excludes the lower decile of the distribution, communicating both the
central tendency and the spread.

\subsection{Interpretation}\label{interpretation}

A reporting sentence: ``The posterior probability had mean 0.71
(posterior \(\mathrm{SD} = 0.11\), 95 \% credible interval 0.47 to
0.89), placing high credibility on values above 0.5 and effectively
excluding the prior-favoured region near 0.3.'' When the posterior is
roughly symmetric the choice of mean versus median rarely matters; when
it is skewed, report both and let readers see the asymmetry.

\subsection{Practical Tips}\label{practical-tips}

\begin{itemize}
\tightlist
\item
  For symmetric posteriors, mean and median are essentially equal;
  choose whichever is conventional in your field.
\item
  For posteriors over bounded parameters (probabilities, variances,
  \(r\)-coefficients), the mode can fall on the boundary; the mean stays
  in the interior and is a safer default.
\item
  Posterior standard deviation is \emph{not} the same as Monte Carlo
  standard error --- the former is genuine uncertainty in \(\theta\),
  the latter is sampling noise from finite MCMC iterations. Report MCSE
  separately when it is appreciable relative to the posterior SD.
\item
  Avoid reporting a ``Bayesian point estimate'' without specifying which
  one; default conventions differ across software and journals.
\item
  For predictions, summarise the posterior predictive distribution
  rather than transforming a point estimate of \(\theta\); only the
  former propagates parameter uncertainty.
\end{itemize}

\subsection{R Packages Used}\label{r-packages-used}

Base R suffices for raw summaries (\texttt{mean}, \texttt{median},
\texttt{var}, \texttt{sd}, \texttt{quantile}),
\texttt{posterior::summarise\_draws()} provides tidy tibbles with Monte
Carlo standard errors, and \texttt{bayestestR::point\_estimate()}
exposes mean / median / MAP through one consistent interface for
\texttt{brms} and \texttt{rstanarm} model objects.

\subsection{For Reviewers}\label{for-reviewers}

\textbf{What to look for in a paper using this method.}

\begin{itemize}
\tightlist
\item
  \textbf{Common misapplications.}
\item
  \textbf{Diagnostics that should be reported but often aren't.}
\item
  \textbf{Red flags in tables and figures.}
\item
  \textbf{What to verify.}
\item
  \textbf{What an adequate Methods paragraph must contain.}
\end{itemize}

\subsection{See also --- labs in this
chapter}\label{see-also-labs-in-this-chapter}

\begin{itemize}
\tightlist
\item
  \href{labs/lab-bayesian-thinking-control-vs-decision-error.Rmd}{Bayesian
  thinking; α control vs decision error}
\item
  \href{labs/lab-brms-stan-loo-hierarchical-models.Rmd}{brms / Stan,
  LOO, hierarchical models}
\end{itemize}

\end{document}
